\documentclass[11pt]{article}
\usepackage{fontspec}
\setmainfont{Times New Roman}
\setsansfont{Arial}
\setmonofont{Consolas}
\usepackage{amsmath,amssymb}
\usepackage{geometry}
\usepackage{microtype}
\geometry{a4paper,margin=3cm}
\setlength{\parindent}{1.25em}
\setlength{\parskip}{0pt}
\emergencystretch=4em
\allowdisplaybreaks
\newcommand{\eps}{\varepsilon}
\newcommand{\pair}[2]{\left(#1\,\middle|\,#2\right)}

\begin{document}

\begin{center}
{\Large\bfseries 3. К инвариантној теорији форм в $n$ пременных\footnote{Доклад, прочитан на Салзбургском собраньу природоизследовательев 1909.}}\par
\vspace{1.2em}
\emph{Jahresbericht d. DMV} 19 (1910), с. 101--104
\end{center}

\vspace{1em}
В пројективној инвариантној теорији форм в $n$ пременных главно пытанье јест решено: конечност система форм јест доказана. Но ряд дальных пытань еще не был размотрен, имено тыцх, кторе односят се к методам изследованьа связи меджу формами. Резултаты, кторе ја добила в односу к таким пытаньам, хочу кратко подати, але ради јасношчи најпрво начртнути поставкы пытань в тернарној области, кде оны сут познане.

Најпрво имам в виду две фундаменталне теоремы симболичного метода: теорему, же все инвариантне творјеньа можно представити симболично, и другу теорему, же все релације, обставајуче меджу инвариантами, достајут се последовательным применьеньем малого числа симболичных идентичностиј; значит, симболичны метод даје все релације без выхода из области инвариантов.\footnote{Штудија: \emph{Methoden zur Theorie der ternären Formen}. II. \S\ 6. (Леипзиг, Teubner, 1889.)} На тој основе потом идут процесы творјеньа форм, симболичны процес свијаньа и несимболичне процесы диференцираньа, теория разложень в ряды и подобно. К тому хочу назвати еще једну теорему, доказану в мојеј дисертацији,\footnote{\emph{Über die Bildung des Formensystems der ternären biquadratischen Form}. \S\ 1. (\emph{Crelle's Journal}, том 134, 1908.)} имено, же все свијаньа можно створити последовательным применьеньем двех можных ``когредиентных'' свијань; под тым при симболичном продукту $s_x t_x u_\sigma u_\tau$ разумеју две свијаньа $(stu)$ и $(\sigma\tau x)$. На тој теореме можно изградити теорију редуцентов и редукције системов форм, аналогично како в бинарној области, како ја показала в мојеј дисертацији.

Ако тепер прејдемо к $n$-арној области, тогда уже прва теорема симболичного метода плати само условно. Познано јест, же уже при кватернарных формах, кроме точковых координат $x$ и равнинских координат $u$, наставајут координаты премых $p_{ik}$, то јест миноры из двех рядов точковых координат. Аналогично в $n$-арној области имамо ряды пременных $p_\rho$, чији поједине елементы сут миноры из $\rho$ рядов $x$, и ряды симболов $S^\rho$, чији елементы понашајут се како миноры из $\rho$ рядов $s$, когредиентных к $u$. Симболично представјенье чрез детерминанты, како познано, плати само тогда, когда ряды симболов $S^\rho$ разкладајут се в ряды $s$ и ты разпределајут се меджу различными детерминантами; реално значенье имајут потом само таке агрегаты детерминантов, кторе зависят од $S^\rho$. Але кторе агрегаты детерминантов то делајут, не можно прегледно установити; тым губят се все остале теоремы, назване в тернарној области.

Тепер хочу дати просто начело, да бы се добило симболично представјенье, експлицитно в рядах $S^\rho$, и с тым такоже остале теоремы. То јест применьенье познаној теоремы о одговарјајучих матрицах: ``Всакому рядку пременных $p_\rho$ взаимно еднозначно одговарја другы ряд $q_{n-\rho}$, чији поједине елементы сут миноры из $n-\rho$ рядов $u$, связаных с $\rho$ рядами $x$ једначеньами $(u\mid x)=0$.'' Одговарјајуче, всакому рядку симболов $S^\rho$ јест приряджен другы $S_{n-\rho}$, кторы понаша се како бы был составльен из $n-\rho$ рядов, когредиентных к $x$.

Та теорема допушча и другу формулацију, угодну за применьеньа: ``Всакому матричному продукту $(v^{(1)}\cdots v^{(\rho)}\mid y^{(1)}\cdots y^{(\rho)})$,\footnote{То јест сума по продуктах одговарјајучих детерминантов, створеных из матрице $\rho$ рядов $v$ и из матрице $\rho$ рядов $y$.} кде ряды другој матрице сут контрагредиентне к рядам првој, приряджује се детерминант; и обратно, всак детерминант можно преобразити в матричны продукт с предписаным числом рядов $\rho$.'' Так детерминант и матричны продукт јавльајут се како полно еквивалентне, и прва теорема симболичного метода добива форму: ``Все инвариантне творјеньа можно симболично представити матричными продуктами.'' Из двех рядов симболов $S^\sigma$, $T^\tau$ с помочју рядов пременных тепер велми лахко можно створити матричны продукт, кторы експлицитно обсахује ты ряды симболов и зато представльа инвариантно творјенье формы $(S^\sigma\mid p_\sigma)(T^\tau\mid p_\tau)$, имено:
\begin{equation}
 \pair{q_{n-\sigma-\lambda}S^\sigma}{T_{n-\tau}p_{\tau-\lambda}},
 \qquad (\lambda=1,2,\ldots,\tau\ \text{або}\ n-\sigma).
\tag{1}
\end{equation}

Важнејша јест обратност: же все инвариантне творјеньа можно експлицитно представити матричными продуктами, так же выше даны процес значи проширенье бинарного и тернарного процеса свијаньа. За доказ тој обратности потребно јест пренести другу теорему симболичного метода на $n$-арну област, имено установити целу множину независных неразложимых идентичностиј, при кторых все ряды $S^\rho$, $p_\rho$ разматривајут се како неразложиме. Те идентичности достајут се из познаных, кторе обсахујут само ряды $x$ и $u$,\footnote{Е. Пасцал в \emph{Memorie d. R. Acc. d. Lincei}. (4) 5 (1888), п. 375.} чрез замену теоремы продукта за детерминанты системом теорем продукта за матрице, и чрез стягненье различных матричных продуктов в једен едины с помочју точно тыцх теорем продукта. Так се доходит до двех дуалистичных формул, из кторых туто наведена јест само једна:
\begin{equation}
\small
\begin{aligned}
 \pair{S_1^{\rho_1}S_2^{\rho_2}\cdots S_k^{\rho_k}}{x p_{\rho-1}}
 &=\sum_{i=1}^{k}\eps\,\pair{S_1^{\rho_1}\cdots S_{i-1}^{\rho_{i-1}}S_{i+1}^{\rho_{i+1}}\cdots S_k^{\rho_k}q_{n-\rho_i+1}}{S_{n-\rho_i}x},\\
 \rho&=\sum_{i=1}^{k}\rho_i<n,\qquad \eps=\pm1.
\end{aligned}
\tag{2}
\end{equation}

Једина специализација, ктора јест в тыцх формулах, имено входенье ряда $x$, не може се објти; при вполно обчих рядах симболов и пременных доходимо до ``разкладној идентичности'', идентичности, в кторој ряд $p_\rho$ разкладаје се на поједине ряды $x$. Найпростейши специалны случай разкладној идентичности уже употребил Clebsch в својој ``фундаменталној задачи инвариантној теорије'' за изводженье елементарных ковариантов при разложеньах в ряды. С помочју тыцх идентичностиј, кторе посредкујут преход од неексплицитных детерминантных выражень к матричному продукту, можно в истине доказати, же с матричным представјеньем добито јест желано експлицитно симболично представјенье.

За процес свијаньа пак плати целком аналогична теорема како в тернарној области; на ньеј можно такоже аналогично изградити редукцијне теоремы. Ако в (1) означимо число $\lambda$ како дефект свијаньа, а свијаньа, за кторе $\sigma=\tau$, како когредиентне, тогда все свијаньа можно створити последовательным применьеньем когредиентных свијань дефекта 1. Доказ тој теоремы основно опира се на факт, же најобчејше формы можно заменити ``нормалными формами'', т. ј. формами, кторе при применьеньу всих инвариантных диференциалных операциј идентично изчезајут. В кватернарној области тут факт доказал Мертенс;\footnote{\emph{Über invariante Gebilde quaternärer Formen}. \emph{Wiener Berichte}. том 98, 1889.} наше познанье најобчејших диференциалных операциј -- кторе, како познано, изходят из процесов свијаньа чрез замену рядов симболов диференциалными симболами -- дозвольаје нам, с применьеньем разкладној идентичности, скоро дословно пренести Мертенсов доказ на $n$-арну област.

\end{document}

